\documentclass[12pt]{article}

% ===================================================================
% Pacotes
% ===================================================================
\usepackage{sbc-template}
\usepackage{graphicx,url}
\usepackage[utf8]{inputenc}      % conteúdo do arquivo em UTF-8
\usepackage[T1]{fontenc}         % acentuação correta na saída
\usepackage[brazil]{babel}
\usepackage{amsmath}
\usepackage{amssymb}
\usepackage{float}
\usepackage{booktabs}            % \toprule, \midrule, \bottomrule
\usepackage{multirow}            % \multirow nas tabelas de resultados

\sloppy

% ===================================================================
% Título e autores  --  PREENCHER os nomes reais dos integrantes
% ===================================================================
\title{Previsão multi-horizonte do preço do ouro:\\
       uma comparação entre modelos lineares, \emph{boosting} e redes
       neurais profundas sobre log-retornos}

\author{Francisco Paulino Arruda Filho \inst{1}, Ivina Lorena Oliveira Moura\inst{1},\\ Gabriel de Souza Mendes\inst{1}, José Vinicius Evangelista Dias de Souza\inst{1}}

\address{
  Mestrado e Doutorado em Ciência da Computação (MDCC)\\
  Universidade Federal do Ceará (UFC) -- Fortaleza, CE -- Brasil
\email{\{paulinofilho, ivinalorena, gabrielmendes01, joseedsouza\}@alu.ufc.br}}


\begin{document}

\maketitle

% ===================================================================
% Resumos
% ===================================================================
\begin{abstract}
This work addresses daily gold price forecasting as a multi-horizon regression problem.
Instead of the raw price, which is highly non-stationary, we adopt as target the
cumulative log-return over $h\!\in\!\{5,15,30\}$ trading days, reconstructed to
US\$/oz for evaluation. From 20 years of data, we build 19 strictly causal,
price‑level‑invariant features and feed 60‑day sliding windows to five models:
linear regression (baseline), Bayesian‑optimized XGBoost, Gated Recurrent Unit(GRU) and Long Short-Term Memory (LSTM) recurrent
networks, and a causal dilated 1D CNN. All share a single temporal split without
leakage and are compared to the \emph{random walk} by the
$\mathrm{MAE}/\mathrm{MAE}_{\text{naive}}$ ratio and by directional accuracy. On a
recent test split, CNN and XGBoost outperform the random walk at all horizons and
reach $\approx\!62\%$ directional accuracy for $h\!=\!30$; under walk‑forward
validation, however, the gain over the baseline is small and strongly dependent on
the market regime.
\end{abstract}

\begin{resumo}
  Este trabalho aborda a previsão diária do preço do ouro como um problema de
  regressão multi-horizonte. Em vez do preço bruto, fortemente
  não-estacionário, adota-se como alvo o log-retorno acumulado em
  $h\!\in\!\{5,15,30\}$ pregões, reconstruído para US\$/oz na avaliação. A
  partir de $20$ anos de dados, construímos $19$ atributos estritamente
  causais e invariantes a nível de preço, e alimentamos janelas deslizantes de
  $60$ dias a cinco modelos: regressão linear (referência), XGBoost com
  otimização bayesiana, redes recorrentes Gated Recurrent Unit(GRU) e Long Short-Term Memory (LSTM) e uma CNN 1D causal com
  dilatação. Todos compartilham um único \emph{split} temporal sem vazamento e
  são comparados ao \emph{random walk} pela razão
  $\mathrm{MAE}/\mathrm{MAE}_{\text{naive}}$ e pela acurácia direcional. Num
  \emph{split} de teste recente, a CNN e o XGBoost superam o \emph{random walk}
  em todos os horizontes e atingem $\approx\!62\%$ de acurácia direcional em
  $h\!=\!30$; sob validação \emph{walk-forward}, contudo, o ganho sobre a
  referência é pequeno e fortemente dependente do regime de mercado.
\end{resumo}

% ===================================================================
\section{Introdução}
\label{sec:introducao}

Séries temporais consistem em sequências de observações coletadas em instantes
sucessivos e, em geral, igualmente espaçados no tempo. A análise dessas séries
busca identificar padrões recorrentes --- tendências, sazonalidades, regimes de
volatilidade --- com o objetivo de apoiar a estimativa de valores futuros
\cite{boxjenkins2015,hyndman2021fpp}. O preço do ouro é um caso particularmente
interessante: trata-se de um ativo de reserva e instrumento clássico de
proteção (\emph{hedge}) contra inflação e incerteza geopolítica, cuja cotação
acompanha de perto crises econômicas, ciclos de juros e a força do dólar.

Prever esse preço, entretanto, é difícil por razões estruturais. Em primeiro
lugar, séries de preços de ativos financeiros aproximam-se de um passeio
aleatório (\emph{random walk}): sob a hipótese de mercados eficientes, a melhor
previsão para o preço de amanhã é, em grande medida, o preço de hoje
\cite{fama1970efficient}. O sinal preditível, quando existe, é fraco e está
imerso em ruído. Em segundo lugar, o nível do preço é fortemente
não-estacionário: como mostra a Figura~\ref{fig:serie}, a cotação do ouro saiu
de cerca de US\$\,280 por onça troy no início dos anos 2000 para mais de
US\$\,3.000 em 2025. Um modelo treinado diretamente no preço bruto jamais
observaria, em treino, os patamares presentes no conjunto de teste --- um
problema de extrapolação que inviabiliza a generalização.

A hipótese de trabalho deste projeto é dupla: (i) reformular o alvo em termos
de \emph{log-retornos acumulados}, que são aproximadamente estacionários, torna
o problema bem-posto; e (ii) modelos capazes de capturar dependências
temporais e relações não-lineares entre indicadores técnicos podem extrair, dos
$60$ dias anteriores, um sinal direcional marginal sobre o \emph{random walk}.
Formalmente, a tarefa é definida como: dada uma janela de $60$ pregões de
atributos derivados do preço, prever o log-retorno acumulado em $5$, $15$ e $30$
dias úteis à frente, simultaneamente (previsão multi-horizonte).

As principais contribuições deste trabalho são:
\begin{itemize}
  \item a definição de um \emph{pipeline} de avaliação rigoroso e sem
        vazamento temporal (\emph{split} cronológico com \emph{gap},
        normalização ajustada apenas no treino, métricas relativas a um
        \emph{baseline} ingênuo), compartilhado por todos os modelos;
  \item a comparação sistemática de cinco famílias de modelos sobre
        \emph{exatamente} as mesmas janelas e o mesmo alvo --- regressão
        linear, \emph{gradient boosting} (XGBoost), redes recorrentes
        (Gated Recurrent Unit(GRU) e Long Short-Term Memory (LSTM)) e rede convolucional 1D causal;
  \item uma análise honesta da robustez via validação \emph{walk-forward},
        que evidencia que o ganho sobre o \emph{random walk} é real, porém
        pequeno e fortemente dependente do regime de mercado.
\end{itemize}

\begin{figure}[H]
  \centering
  \includegraphics[width=0.95\linewidth]{figuras/serie_preco.png}
  \caption{Preço de fechamento diário do ouro (US\$/oz) entre 2000 e 2025. A
    forte tendência de alta e a ausência de média constante caracterizam uma
    série não-estacionária, motivando o uso do log-retorno como alvo.}
  \label{fig:serie}
\end{figure}

% ===================================================================
\section{Fundamentação Teórica}
\label{sec:fundamentacao}

\subsection{Definição do problema}
\label{sec:definicao}

Formaliza-se a tarefa como um problema de regressão multi-horizonte sobre uma
série temporal financeira. Seja $P_t$ o preço de fechamento do ouro no pregão
$t$. A cada instante dispõe-se de uma janela deslizante de $L = 60$ pregões
contendo um vetor de atributos $\mathbf{x}_\tau \in \mathbb{R}^{d}$ por dia,
formando a matriz de entrada
\begin{equation}
  X_t = [\mathbf{x}_{t-L+1}, \dots, \mathbf{x}_t] \in \mathbb{R}^{L \times d}.
\end{equation}
O objetivo não é prever o preço bruto, mas o \emph{log-retorno acumulado} em
$h$ dias à frente:
\begin{equation}
  y_t^{(h)} = \log\!\left(\frac{P_{t+h}}{P_t}\right),
  \qquad h \in \{5, 15, 30\}.
  \label{eq:alvo}
\end{equation}
A escolha do log-retorno como alvo justifica-se por dois motivos. Primeiro, como
discutido na Seção~\ref{sec:introducao}, o preço é fortemente não-estacionário
(ver Figura~\ref{fig:serie}), enquanto o log-retorno é aproximadamente
estacionário e centrado em zero, evitando o problema de extrapolação. Segundo,
o log-retorno é aditivo no tempo e simétrico, o que facilita tanto a
modelagem quanto a reconstrução do preço. Esta última é feita
\emph{a posteriori} por
\begin{equation}
  \hat{P}_{t+h} = P_t \cdot \exp\!\big(\hat{y}_t^{(h)}\big),
  \label{eq:reconstrucao}
\end{equation}
de modo que todas as métricas de erro possam ser expressas em US\$/oz,
interpretáveis pelo usuário final.

\subsection{Descrição dos dados}
\label{sec:dados}

Os dados provêm do conjunto público \emph{Precious Metals History since 2000
with News}, disponível no Kaggle, com cotações diárias de metais preciosos
desde 2000 e manchetes associadas às datas. Utilizou-se a série do ouro,
cobrindo mais de duas décadas de pregões (de 2000 a 2025). As colunas
originais relevantes são:

\begin{description}
  \item[\texttt{timestamp}] data da observação, com granularidade diária;
        índice temporal de todos os \emph{merges} e janelas móveis.
  \item[\texttt{open, high, low, close}] preços de abertura, máxima, mínima e
        fechamento do dia, na unidade de cotação.
  \item[\texttt{volume}] volume negociado do ativo no dia.
  \item[\texttt{currency, unit}] moeda (USD) e unidade (onça troy) de
        denominação --- campos categóricos, constantes em toda a série e,
        portanto, descartados na modelagem.
  \item[\texttt{headlines}] texto de manchete associado à data, de natureza
        textual.
\end{description}

\subsubsection{Engenharia de atributos}
\label{sec:features}

A partir da série de preços derivaram-se $d = 19$ atributos exclusivamente
\emph{relativos} (retornos, razões e osciladores), nunca o preço bruto, de modo
a generalizar entre diferentes níveis de cotação. O conjunto, descrito na
Tabela~\ref{tab:features}, reúne log-retornos em múltiplas escalas,
volatilidades realizadas e sua razão (indicador de estresse), distância a
médias móveis, cruzamento de médias, Relative Strength Index (RSI), histograma Moving Average Convergence Divergence (MACD), posição nas bandas
de Bollinger, \emph{drawdown} de $252$ dias, descritores intradiários da barra
Open–High–Low–Close (OHLC), e um índice textual de ``medo'' extraído das manchetes.

O índice de medo é calculado por um léxico de termos associados a crises
geopolíticas e financeiras (por exemplo, \emph{crisis}, \emph{crash},
\emph{war}, \emph{recession}, \emph{inflation}, \emph{default},
\emph{pandemic}), contando ocorrências por manchete e suavizando por média
móvel de $5$ dias. Trata-se de uma abordagem deliberadamente simples de Processamento de Linguagem Natural (PLN)
baseada em contagem; alternativas mais sofisticadas (\emph{embeddings},
modelos de sentimento financeiro) ficam como trabalho futuro.

\begin{table}[H]
  \caption{Os $19$ atributos de entrada. Todos são causais (usam apenas
    informação até $t$) e invariantes ao nível de preço.}
  \label{tab:features}
  \centering
  \small
  \begin{tabular}{lll}
    \toprule
    Atributo & Definição & Grupo \\
    \midrule
    \texttt{log\_ret\_1,5,21} & $\log(P_t/P_{t-k})$, $k\in\{1,5,21\}$ & Momentum \\
    \texttt{vol\_5, vol\_21}  & desvio-padrão móvel de \texttt{log\_ret\_1} & Volatilidade \\
    \texttt{vol\_ratio}       & $\mathrm{vol}_5/\mathrm{vol}_{21}$ ($>1$ = estresse) & Volatilidade \\
    \texttt{dist\_ma20,50}    & $P_t/\mathrm{MA}_k - 1$, $k\in\{20,50\}$ & Tendência \\
    \texttt{ma\_cross}        & $\mathrm{MA}_{20}/\mathrm{MA}_{50} - 1$ & Tendência \\
    \texttt{rsi}              & RSI(14) recentrado em $[-1,1]$ & Oscilador \\
    \texttt{macd\_hist}       & histograma MACD normalizado por $P_t$ & Oscilador \\
    \texttt{boll\_b}          & posição \%B nas bandas de Bollinger & Oscilador \\
    \texttt{drawdown\_252}    & $P_t/\max_{252}(P) - 1$ & Regime \\
    \texttt{hl\_range}        & $(\mathrm{high}-\mathrm{low})/P_t$ & Intradiário \\
    \texttt{close\_pos}       & posição do fechamento na barra OHLC & Intradiário \\
    \texttt{gap\_open}        & $\log(\mathrm{open}_t/P_{t-1})$ & Intradiário \\
    \texttt{volume}           & volume negociado no dia & Liquidez \\
    \texttt{fear\_ratio}      & índice de medo por léxico nas manchetes & Texto \\
    \texttt{fear\_ma5}        & média móvel de $5$ dias de \texttt{fear\_ratio} & Texto \\
    \bottomrule
  \end{tabular}
\end{table}

A análise exploratória dos dados (Seção~\ref{sec:eda}) confirma duas
propriedades estatísticas centrais: o preço é não-estacionário (o teste de
Dickey--Fuller Aumentado, ADF, não rejeita a hipótese de raiz unitária no
preço, mas a rejeita nos log-retornos) \cite{dickeyfuller1979}, e a
distribuição dos log-retornos é leptocúrtica, com caudas pesadas e assimetria
negativa --- justificando o uso de funções de perda e normalizadores robustos
(Seção~\ref{sec:preproc}).

\subsection{Trabalhos relacionados}
\label{sec:relacionados}

A previsão de séries financeiras com aprendizado profundo é um campo ativo com trabalhos relevantes.
\cite{sezer2020dlreview} revisam sistematicamente o uso de redes neurais para
esse fim entre 2005 e 2019, destacando o predomínio de redes recorrentes
(LSTM/GRU) e o crescente interesse por arquiteturas convolucionais e de atenção.
As redes LSTM \cite{hochreiter1997lstm} e GRU \cite{cho2014gru} são padrão para
dependências temporais de longo prazo. Em paralelo, modelos convolucionais
causais, de \emph{WaveNet} \cite{oord2016wavenet} às \emph{Temporal
Convolutional Networks} (TCN) \cite{bai2018tcn}, mostraram que convoluções
dilatadas competem com (ou superam) (Recurrent Neural Networks) RNNs em modelagem de sequências, com
treinamento mais estável e paralelizável. Para dados tabulares derivados de
janelas, métodos de \emph{gradient boosting} como o XGBoost
\cite{chen2016xgboost} são fortíssimos \emph{baselines}. Mais recentemente, o
\emph{Temporal Fusion Transformer} \cite{lim2021tft} unificou previsão
multi-horizonte e interpretabilidade. Embora não tenha sido implementado aqui
por restrições de tempo, é uma direção natural de continuidade.

O ponto de partida direto deste trabalho foi um experimento acadêmico anterior
(2025) que utilizou \emph{este mesmo conjunto de dados} com \emph{Ridge
Regression} para prever o ouro um dia à frente. Aquele trabalho continha uma
falha metodológica grave: aplicou \emph{shuffle} aleatório nas amostras antes do
\emph{split}, misturando passado e futuro e produzindo vazamento de dados e
\emph{overfitting} otimista. Código disponível em
\url{https://github.com/JoseEdSouza/machine-learning-works/blob/main/gold/eda.ipynb}.
Corrigir esse erro, adotando um protocolo estritamente temporal com
\emph{gap} anti-vazamento, é uma preocupação central da
Seção~\ref{sec:metodologia}.

% ===================================================================
\section{Metodologia}
\label{sec:metodologia}

\subsection{Pré-processamento e engenharia de atributos}
\label{sec:preproc}

Todos os atributos da Tabela~\ref{tab:features} são \emph{causais}: em cada
instante $t$ utilizam apenas informação disponível até $t$, não havendo
vazamento temporal a partir das \emph{features}. As janelas deslizantes de
$L = 60$ dias foram construídas descartando-se janelas que contivessem qualquer
valor ausente (em geral no início da série, por conta das médias móveis longas).

A padronização foi feita com \texttt{RobustScaler} (mediana e intervalo
interquartil), mais adequado às caudas pesadas típicas de períodos de crise do
que a padronização por média/desvio, seguido de \emph{clipping} em $\pm 8$
desvios robustos para limitar a influência de dias de pânico sem descartá-los.
Crucialmente, o normalizador foi ajustado \emph{exclusivamente} sobre o conjunto
de treino e então aplicado a validação e teste, evitando o vazamento de
estatísticas do futuro para o passado. O mesmo \texttt{RobustScaler} é aplicado
ao alvo (log-retornos) para facilitar a convergência das redes; as predições são
revertidas à escala original antes do cálculo das métricas.

\subsection{Protocolo de avaliação}
\label{sec:protocolo}

Adotou-se um \emph{split} estritamente temporal (sem embaralhamento),
reservando $70\%$ dos pregões para treino, $15\%$ para validação e $15\%$ para
teste, em ordem cronológica. Para impedir que janelas vizinhas compartilhassem
dias entre conjuntos --- ou que o horizonte do alvo de um conjunto invadisse as
\emph{features} do conjunto seguinte --- inseriu-se um \emph{gap} de
$L + \max(h) = 60 + 30 = 90$ pregões entre conjuntos. Com esse espaçamento, o
último alvo de cada conjunto encerra antes do início da primeira janela de
\emph{features} do conjunto seguinte, garantindo ausência de sobreposição. Como
ilustra a Figura~\ref{fig:split}, essa divisão posiciona a crise de 2008 no
treino, a pandemia de 2020 na validação e o forte \emph{rally} de 2021--2025 no
teste --- um cenário desafiador, pois exige extrapolar para níveis de preço
nunca vistos (viabilizado justamente pelo alvo em log-retorno).

\begin{figure}[tbp]
  \centering
  \includegraphics[width=0.95\linewidth]{figuras/split_temporal.png}
  \caption{Divisão temporal treino+validação / teste com \emph{gap}
    anti-vazamento. O período de teste corresponde ao \emph{rally} recente do
    ouro (2021--2025), exigindo generalização para patamares de preço inéditos.}
  \label{fig:split}
\end{figure}

Para uma avaliação mais robusta dos modelos neurais, complementou-se o
\emph{split} único com validação \emph{walk-forward} (janela de treino
expansiva, fold de teste de $252$ pregões $\approx 1$ ano, com embargo de
$L + \max(h)$ amostras). Como discutido por \cite{bergmeir2012cv}, a validação
cruzada com respeito à ordem temporal é apropriada para séries; a
\emph{walk-forward} fornece, além disso, uma distribuição de desempenho ao
longo de diferentes anos/regimes, e não um único número possivelmente
afortunado.

\subsection{Modelos avaliados}
\label{sec:modelos}

A estratégia de modelagem segue uma progressão de complexidade. A regressão
linear, adotada como linha de base, fornece um piso de desempenho a partir do
qual se mede o ganho efetivo de arquiteturas mais sofisticadas; estas só se
justificam quando superam significativamente essa fundação mínima.

\subsubsection{Referências: \emph{random walk} e regressão linear}
\label{sec:baseline}

A referência inferior é o \emph{baseline} ingênuo de passeio aleatório
(\emph{random walk}): o melhor palpite para o preço em $t+h$ é o preço em $t$,
isto é, log-retorno previsto igual a zero ($\hat{P}_{t+h} = P_t$). Qualquer
modelo que não supere esse \emph{baseline} não aprendeu nada útil.

Como \emph{baseline} treinável, adotou-se a regressão linear múltipla,
ajustada por mínimos quadrados ordinários
($\hat{\boldsymbol{\beta}} = (\mathbf{X}^\top \mathbf{X})^{-1}\mathbf{X}^\top \mathbf{y}$),
da biblioteca \texttt{scikit-learn} \cite{pedregosa2011sklearn}. Para colocá-la
em pé de igualdade com os demais modelos, usou-se como entrada a \emph{mesma}
janela deslizante de $L \times d$ atributos, achatada em um vetor de
$60 \times 19 = 1140$ regressores. A função de autocorrelação parcial (PACF) dos
log-retornos foi usada na análise exploratória para verificar a (quase) ausência
de autocorrelação linear significativa.

\subsubsection{XGBoost com otimização bayesiana}
\label{sec:xgboost}

O primeiro modelo não-linear é o XGBoost \cite{chen2016xgboost}, um
\emph{gradient boosting} de árvores, treinado de forma independente por
horizonte sobre a janela achatada de atributos. Os hiperparâmetros foram
otimizados por busca bayesiana (\texttt{BayesSearchCV}, \texttt{scikit-optimize})
\cite{snoek2012bayesopt} com $16$ iterações por horizonte e validação cruzada
temporal de $5$ \emph{folds} (\texttt{TimeSeriesSplit}). O espaço de busca cobriu
profundidade ($3$--$5$), taxa de aprendizado ($0{,}02$--$0{,}10$, log-uniforme),
\emph{subsample} e \emph{colsample} ($0{,}75$--$0{,}95$), número de árvores
($200$--$500$) e regularização $L_1/L_2$.

\subsubsection{Redes recorrentes: GRU e LSTM}
\label{sec:rnn}

Avaliaram-se duas redes recorrentes que processam a janela de $60$ dias passo a
passo: a GRU \cite{cho2014gru} e a LSTM \cite{hochreiter1997lstm}, esta última
com uma célula de memória explícita ao custo de mais parâmetros. As duas
arquiteturas são idênticas a menos da célula recorrente: duas camadas
recorrentes de $64$ unidades intercaladas por \emph{Layer Normalization}
\cite{ba2016layernorm} e \emph{Spatial Dropout} \cite{tompson2015spatialdropout},
seguidas de uma camada densa de $64$ unidades com \emph{dropout}
\cite{srivastava2014dropout} e \textbf{três cabeças de saída} (uma por
horizonte) que compartilham a mesma representação, numa formulação de
aprendizado multitarefa. A perda é a de \emph{Huber} \cite{huber1964}, robusta a
\emph{outliers} de retorno, ponderada por horizonte ($1{,}0$, $0{,}7$, $0{,}5$),
com otimizador Adam \cite{kingma2015adam}.

\subsubsection{Rede convolucional 1D causal (modelo principal)}
\label{sec:cnn}

O modelo \emph{champion} é uma rede convolucional unidimensional (CNN 1D) com
convoluções \emph{causais} --- que não acessam instantes futuros dentro da
janela --- e dilatação crescente, à maneira de \emph{WaveNet}/TCN
\cite{oord2016wavenet,bai2018tcn}. Empregaram-se três camadas convolucionais
com dilatação $1 \to 2 \to 4$, o que amplia exponencialmente o campo receptivo
e cobre os $60$ dias com poucas camadas, cada uma seguida de \emph{Layer
Normalization}. Após \emph{Spatial Dropout}, a representação temporal é resumida
pela concatenação de \emph{pooling} médio e máximo globais, capturando
simultaneamente tendência (média) e choques abruptos (máximo). Seguem uma camada
densa de $64$ unidades com \emph{dropout} e, novamente, três cabeças de saída
$h \in \{5,15,30\}$. A perda é a de \emph{Huber}; o otimizador é o Adam. A
arquitetura completa tem cerca de $60$ mil parâmetros treináveis. O treinamento
usou \emph{early stopping} com restauração dos melhores pesos e redução da taxa
de aprendizado em platôs, ambos monitorando a perda de validação.

\subsubsection{Extensão: \emph{walk-forward} com previsão por quantis}
\label{sec:quantis}

Como evolução do modelo pontual, treinou-se uma variante da CNN que prevê, para
cada horizonte, três quantis ($q_{10}, q_{50}, q_{90}$) em vez de um único
valor, usando a \emph{pinball loss} (perda quantílica)
\cite{koenkerbassett1978}. Combinada à validação \emph{walk-forward}, essa
variante produz intervalos de incerteza ($p10$--$p90$) por horizonte, tornando a
previsão mais informativa na prática do que uma estimativa pontual.

\subsection{Métricas}
\label{sec:metricas}

Os modelos foram avaliados por Root Mean Squared Error (RMSE), Mean Absolute Error (MAE) e Mean Absolute Percentage Error (MAPE), calculados sobre o
\emph{preço reconstruído} (Eq.~\ref{eq:reconstrucao}), e pela \emph{acurácia
direcional} (proporção de acertos quanto ao sinal do retorno --- a métrica de
maior relevância prática para decisão de compra/venda). Como referência
sintética, reporta-se a razão
\begin{equation}
  \frac{\mathrm{MAE}}{\mathrm{MAE}_{\text{naive}}},
\end{equation}
em que o \emph{baseline} ingênuo (\emph{random walk}) prevê preço constante
$\hat{P}_{t+h} = P_t$. Em palavras, a razão mede o erro do modelo em proporção
ao de um previsor ingênuo que apenas repete o último preço observado: valores
abaixo de $1$ indicam ganho sobre essa referência (por exemplo, $0{,}95$ = erro
$5\%$ menor que o \emph{random walk}), $1{,}0$ indica empate, e valores acima
de $1$ indicam desempenho \emph{pior} do que não prever nada.

Ressalta-se que métricas calculadas diretamente na escala de preço, como o
$R^2$, são \emph{enganosas} neste problema: por serem dominadas pelo nível
não-estacionário do preço, atribuem desempenho artificialmente alto mesmo a
modelos triviais --- um ponto que será evidenciado empiricamente na
Seção~\ref{sec:discussao}. Por isso a avaliação privilegia métricas relativas ao
\emph{baseline} e na escala de retorno. Para a vertente de volatilidade,
reportam-se ainda a perda \emph{QLIKE}, comum em finanças por penalizar
sub-previsão de risco \cite{patton2011volatility,zhang2015rsgarch}, e o
coeficiente de correlação de postos de Spearman ($\rho$), medida não-paramétrica
da relação monotônica entre previsto e realizado \cite{dellinger2017spearman}.

% ===================================================================
\section{Experimentos e Resultados}
\label{sec:experimentos}

\subsection{Configuração experimental}
\label{sec:config}

Os experimentos foram implementados em Python, com \texttt{scikit-learn}
\cite{pedregosa2011sklearn} e \texttt{XGBoost} \cite{chen2016xgboost} para os
modelos clássicos e \texttt{TensorFlow}/\texttt{Keras}
\cite{abadi2016tensorflow,chollet2015keras} para as redes neurais, com
aceleração em GPU. A reprodutibilidade foi garantida fixando a semente
aleatória ($\texttt{SEED}=42$). Os hiperparâmetros das redes neurais foram
selecionados com \texttt{Keras Tuner}; os do XGBoost, com busca bayesiana. Os
artefatos (métricas, modelos, gráficos) são versionados por experimento, e o
treinamento foi monitorado com \texttt{TensorBoard}. A
Tabela~\ref{tab:hiperparametros} resume a configuração final do modelo
principal e das redes recorrentes.

\begin{table}[H]
  \caption{Configuração final dos modelos neurais (selecionada por busca de
    hiperparâmetros). Janela $L=60$, horizontes $\{5,15,30\}$, perda de Huber,
    otimizador Adam.}
  \label{tab:hiperparametros}
  \centering
  \small
  \begin{tabular}{lll}
    \toprule
    Hiperparâmetro & CNN 1D (principal) & GRU / LSTM \\
    \midrule
    Camadas extratoras & 3 $\times$ Conv1D causal & 2 $\times$ recorrente (64) \\
    Filtros / unidades & 96, 64, 256 & 64, 64 \\
    Dilatação / kernel & $1{\to}2{\to}4$ / 5,3,3 & --- \\
    \emph{Spatial Dropout} & 0{,}20 & 0{,}15 \\
    Densa + \emph{Dropout} & 64 / 0{,}20 & 64 / 0{,}25 \\
    Taxa de aprendizado & $4{,}1\times10^{-3}$ & $1{,}0\times10^{-3}$ \\
    Pesos de perda ($h$5,15,30) & 1{,}0 / 0{,}6 / 0{,}2 & 1{,}0 / 0{,}7 / 0{,}5 \\
    Épocas / \emph{batch} & 120 / 64 & 120 / 64 \\
    \emph{Early stopping} & paciência 15 & paciência 15 \\
    Parâmetros treináveis & $\approx 60$ mil & $\approx 60$--$80$ mil \\
    \bottomrule
  \end{tabular}
\end{table}

\subsection{Análise exploratória dos dados}
\label{sec:eda}

A análise exploratória confirma as dificuldades estruturais do problema. A
distribuição dos log-retornos (Figura~\ref{fig:retornos}) é claramente
leptocúrtica: mais concentrada em torno de zero e com caudas muito mais pesadas
que a normal teórica, além de assimetria negativa (quedas extremas mais
frequentes que altas extremas). Esse comportamento motiva tanto o
\texttt{RobustScaler} quanto a perda de Huber. A
Figura~\ref{fig:acfpacf} mostra Autocorrelation Function (ACF) e PACF dos log-retornos: a autocorrelação
linear é fraca e, em sua maior parte, dentro das bandas de confiança ---
evidência direta da proximidade com um passeio aleatório e da insuficiência de
modelos puramente lineares. Por fim, a Figura~\ref{fig:vol} exibe a volatilidade
anualizada (janela de $21$ dias), revelando forte agrupamento de volatilidade
(\emph{volatility clustering}), com picos nas crises de 2008 e 2020 ---
fenômeno classicamente modelado por famílias (Generalized Autoregressive Conditional Heteroskedasticity (GARCH) \cite{bollerslev1986garch}.

\begin{figure}[tbp]
  \centering
  \includegraphics[width=0.85\linewidth]{figuras/distribuicao_retornos.png}
  \caption{Distribuição dos log-retornos diários: histograma \emph{vs.} normal
    teórica, Q-Q plot e \emph{boxplot}. As caudas pesadas e a assimetria
    negativa evidenciam leptocurtose.}
  \label{fig:retornos}
\end{figure}

\begin{figure}[H]
  \centering
  \includegraphics[width=0.85\linewidth]{figuras/acf_pacf.png}
  \caption{ACF e PACF dos log-retornos. A autocorrelação linear é fraca e em
    grande parte não-significativa, coerente com a hipótese de passeio
    aleatório.}
  \label{fig:acfpacf}
\end{figure}

\begin{figure}[H]
  \centering
  \includegraphics[width=0.95\linewidth]{figuras/volatilidade.png}
  \caption{Volatilidade anualizada (janela móvel de $21$ dias). Os picos
    coincidem com as crises de 2008 e 2020, em destaque, ilustrando o agrupamento de
    volatilidade.}
  \label{fig:vol}
\end{figure}

\subsection{Resultados no conjunto de teste}
\label{sec:resultados}

A Tabela~\ref{tab:resultados} consolida o desempenho de todos os modelos no
conjunto de teste, por horizonte. Três achados se destacam.

\textbf{(i) A regressão linear não bate o \emph{random walk}.} Apesar de
reduzir o erro nos horizontes mais longos, o \emph{baseline} linear apresenta
$\mathrm{MAE}/\mathrm{MAE}_{\text{naive}} > 1$ em $h=5$ e $h=15$ ($1{,}160$ e
$1{,}068$): com $1140$ regressores e sinal fraco, o modelo linear sobre-ajusta
e não generaliza melhor que prever ``preço constante''. Isso confirma que
estruturas lineares são insuficientes e estabelece o piso que os demais modelos
precisam superar.

\textbf{(ii) CNN e XGBoost superam o \emph{random walk} em todos os
horizontes.} A CNN 1D causal obtém a melhor razão $\mathrm{MAE}/\mathrm{naive}$
em $h=5$ ($0{,}983$) e $h=15$ ($0{,}971$), enquanto o XGBoost lidera em $h=30$
($0{,}951$); ambos atingem $\approx 62\%$ de acurácia direcional em $h=30$,
bem acima dos $50\%$ de uma moeda justa. As redes recorrentes (GRU e LSTM) ficam
ligeiramente atrás, com ganhos marginais sobre a referência.

\textbf{(iii) O ganho é maior em horizontes longos.} Para todos os modelos, a
razão $\mathrm{MAE}/\mathrm{naive}$ cai e a acurácia direcional sobe à medida que
$h$ cresce. Isso é intuitivo: em horizontes curtos o ruído domina e o
\emph{random walk} é difícil de bater; em horizontes mais longos, componentes de
tendência tornam-se relativamente mais previsíveis.

\begin{table}[H]
  \caption{Desempenho no conjunto de teste (\emph{split} único cronológico),
    por horizonte $h$. Erros em US\$/oz sobre o preço reconstruído. A melhor
    razão $\mathrm{MAE}/\mathrm{naive}$ de cada bloco está em \textbf{negrito}.
    A referência \emph{naive} (\emph{random walk}) tem razão $1{,}000$ por
    construção.}
  \label{tab:resultados}
  \centering
  \small
  \begin{tabular}{llccccc}
    \toprule
    $h$ & Modelo & MAE & RMSE & MAPE (\%) & Acur.\ dir.\ (\%) & MAE/naive \\
    \midrule
    \multirow{5}{*}{5}
      & Regressão Linear (\emph{baseline}) & 38{,}23 & 48{,}20 & 1{,}82 & 52{,}5 & 1{,}160 \\
      & XGBoost                            & 32{,}53 & 41{,}62 & 1{,}56 & 56{,}5 & 0{,}987 \\
      & GRU                                & 32{,}65 & 41{,}77 & 1{,}57 & 53{,}4 & 0{,}990 \\
      & LSTM                               & 32{,}62 & 41{,}50 & 1{,}57 & 53{,}3 & 0{,}989 \\
      & \textbf{CNN 1D (principal)}        & 32{,}39 & 41{,}42 & 1{,}56 & 57{,}4 & \textbf{0{,}983} \\
    \midrule
    \multirow{5}{*}{15}
      & Regressão Linear (\emph{baseline}) & 63{,}87 & 80{,}22 & 3{,}05 & 55{,}0 & 1{,}068 \\
      & XGBoost                            & 58{,}26 & 73{,}11 & 2{,}78 & 58{,}9 & 0{,}974 \\
      & GRU                                & 59{,}56 & 74{,}36 & 2{,}84 & 56{,}5 & 0{,}996 \\
      & LSTM                               & 58{,}48 & 73{,}23 & 2{,}78 & 54{,}8 & 0{,}978 \\
      & \textbf{CNN 1D (principal)}        & 58{,}08 & 72{,}68 & 2{,}77 & 58{,}9 & \textbf{0{,}971} \\
    \midrule
    \multirow{5}{*}{30}
      & Regressão Linear (\emph{baseline}) & 87{,}66 & 109{,}09 & 4{,}16 & 58{,}4 & 1{,}002 \\
      & \textbf{XGBoost}                   & 83{,}16 & 101{,}45 & 3{,}89 & 62{,}0 & \textbf{0{,}951} \\
      & GRU                                & 85{,}00 & 103{,}70 & 3{,}97 & 58{,}7 & 0{,}972 \\
      & LSTM                               & 85{,}77 & 103{,}49 & 4{,}02 & 57{,}6 & 0{,}980 \\
      & CNN 1D (principal)                 & 83{,}43 & 102{,}38 & 3{,}92 & 62{,}0 & 0{,}954 \\
    \bottomrule
  \end{tabular}
\end{table}

A Figura~\ref{fig:curvas} mostra as curvas de aprendizado do modelo principal: a
perda de treino cai continuamente, mas a de validação atinge um mínimo por volta
da época $10$--$15$ e volta a subir --- daí a importância do \emph{early
stopping} com restauração dos melhores pesos. A Figura~\ref{fig:pred} compara,
no conjunto de teste, o preço real e o previsto pela CNN nos três horizontes: a
aderência é boa em $h=5$ e o erro/atraso cresce em $h=15$ e $h=30$, como
esperado, com o modelo seguindo a tendência mas suavizando os movimentos mais
bruscos.

\begin{figure}[H]
  \centering
  \includegraphics[width=0.95\linewidth]{figuras/cnn_curvas_treino.png}
  \caption{Curvas de aprendizado da CNN principal. À esquerda, perda total de
    treino e validação; à direita, perda de validação por horizonte. A subida da
    perda de validação após $\approx$ 10 épocas motiva o \emph{early stopping}.}
  \label{fig:curvas}
\end{figure}

\begin{figure}[H]
  \centering
  \includegraphics[width=0.82\linewidth]{figuras/cnn_pred_vs_real.png}
  \caption{Preço real \emph{vs.} previsto pela CNN principal no conjunto de
    teste, para $h=5$, $15$ e $30$ dias. A previsão acompanha a tendência e
    suaviza os choques; o atraso aumenta com o horizonte.}
  \label{fig:pred}
\end{figure}

\subsection{Validação cruzada e dependência de regime}
\label{sec:robustez}

Os números da Tabela~\ref{tab:resultados} referem-se a um único \emph{split} de
teste --- favorável, por coincidir com um forte \emph{rally}. Para uma leitura
honesta, a Tabela~\ref{tab:cv} resume a validação \emph{walk-forward} dos
modelos neurais ao longo de múltiplos anos. O quadro muda: a razão
$\mathrm{MAE}/\mathrm{naive}$ média fica \emph{próxima de} $1{,}0$ (entre
$0{,}99$ e $1{,}07$) com \emph{desvio-padrão alto}, especialmente em $h=30$.
A Figura~\ref{fig:cv} confirma visualmente: as caixas de
$\mathrm{MAE}/\mathrm{naive}$ cruzam a linha de $1{,}0$, e o erro cresce com o
horizonte.

\begin{table}[H]
  \caption{Validação \emph{walk-forward} (folds anuais) dos modelos neurais:
    média $\pm$ desvio-padrão entre folds. A vantagem sobre o \emph{random walk}
    é pequena e instável.}
  \label{tab:cv}
  \centering
  \small
  \begin{tabular}{llcc}
    \toprule
    Modelo & $h$ & Acur.\ dir.\ (\%) & MAE/naive \\
    \midrule
    \multirow{3}{*}{CNN 1D}
      & 5  & $\mathbf{55{,}3 \pm 5{,}7}$  & $\mathbf{0{,}995 \pm 0{,}024}$ \\
      & 15 & $54{,}8 \pm 8{,}4$            & $\mathbf{1{,}006 \pm 0{,}055}$ \\
      & 30 & $54{,}8 \pm 13{,}9$           & $\mathbf{1{,}037 \pm 0{,}138}$ \\
    \midrule
    \multirow{3}{*}{GRU}
      & 5  & $51{,}5 \pm 2{,}0$  & $1{,}042 \pm 0{,}038$ \\
      & 15 & $53{,}9 \pm 6{,}6$  & $1{,}047 \pm 0{,}067$ \\
      & 30 & $\mathbf{58{,}1 \pm 12{,}1}$ & $1{,}074 \pm 0{,}225$ \\
    \midrule
    \multirow{3}{*}{LSTM}
      & 5  & $51{,}4 \pm 4{,}0$  & $1{,}048 \pm 0{,}047$ \\
      & 15 & $\mathbf{55{,}1 \pm 4{,}5}$  & $1{,}049 \pm 0{,}068$ \\
      & 30 & $55{,}9 \pm 10{,}7$ & $1{,}075 \pm 0{,}193$ \\
    \bottomrule
  \end{tabular}
\end{table}

\begin{figure}[H]
  \centering
  \includegraphics[width=0.98\linewidth]{figuras/cnn_cv_boxplot.png}
  \caption{Distribuição das métricas da CNN entre folds da validação
    \emph{walk-forward}, por horizonte (1=$h$5, 2=$h$15, 3=$h$30). As caixas de
    $\mathrm{MAE}/\mathrm{naive}$ cruzam $1{,}0$: o ganho sobre o \emph{random
    walk} não é consistente entre regimes.}
  \label{fig:cv}
\end{figure}

A causa fica clara ao decompor o desempenho por ano (Figura~\ref{fig:anual},
para $h=30$): em anos de tendência definida (2021, 2024 e, sobretudo, 2025) a
acurácia direcional ultrapassa $80\%$ e a $\mathrm{MAE}/\mathrm{naive}$ cai
abaixo de $0{,}85$; já em anos de mercado lateral e ruidoso (2022 e 2023) a
acurácia recua para $\approx 46$--$48\%$ (pior que o acaso) e a razão de erro
ultrapassa $1{,}0$. Em outras palavras, os modelos capturam tendência quando ela existe, mas não ``criam'' sinal num mercado sem direção --- exatamente o que a teoria de
mercados eficientes prevê \cite{fama1970efficient,jegadeesh1993momentum}.

\begin{figure}[H]
  \centering
  \includegraphics[width=0.95\linewidth]{figuras/cnn_anual_h30.png}
  \caption{Desempenho da CNN por ano calendário em $h=30$: acurácia direcional
    (esq.) e $\mathrm{MAE}/\mathrm{naive}$ (dir.). O ganho concentra-se em anos
    de tendência (2021, 2024, 2025) e desaparece em mercados laterais (2022,
    2023).}
  \label{fig:anual}
\end{figure}

\subsection{Discussão}
\label{sec:discussao}

\paragraph{O $R^2$ na escala de preço é enganoso.} A
Figura~\ref{fig:scatter} reporta o $R^2$ entre preço previsto e real:
$0{,}984$, $0{,}954$ e $0{,}920$ para $h=5,15,30$. Esses valores
\emph{altíssimos} dariam a falsa impressão de um modelo quase perfeito ---
mas o \emph{mesmo} modelo, medido pela $\mathrm{MAE}/\mathrm{naive}$, apenas
empata ou supera marginalmente o passeio aleatório. A explicação é que o $R^2$
no preço é dominado pelo nível/tendência da série: até prever ``preço de hoje''
produziria $R^2$ elevado. Esse contraste é a justificativa empírica para
priorizarmos métricas relativas ao \emph{baseline} e na escala de retorno
(Seção~\ref{sec:metricas}).

\begin{figure}[H]
  \centering
  \includegraphics[width=0.98\linewidth]{figuras/cnn_scatter.png}
  \caption{Dispersão preço previsto \emph{vs.} real (CNN, teste). O $R^2$ na
    escala de preço é muito alto ($0{,}92$--$0{,}98$), porém enganoso: ele é
    inflado pelo nível da série e não reflete o ganho real sobre o \emph{random
    walk}.}
  \label{fig:scatter}
\end{figure}

\paragraph{Convoluções, recorrência ou \emph{boosting}?} No \emph{split}
único, a CNN 1D causal e o XGBoost foram os melhores e praticamente empataram,
superando GRU e LSTM. A competitividade das convoluções causais frente às RNNs é
coerente com a literatura de modelagem de sequências
\cite{bai2018tcn,oord2016wavenet}; a força do XGBoost reforça que um
\emph{baseline} não-neural bem ajustado é difícil de bater em dados tabulares.
Na validação \emph{walk-forward}, a CNN apresentou a $\mathrm{MAE}/\mathrm{naive}$
média mais próxima de $1$ (e menos instável), o que a torna a escolha de
modelo principal.

\paragraph{Onde os modelos falham.} O ponto fraco é sistemático: mercados
laterais (2022--2023). Nesses regimes não há tendência a ser capturada e o sinal
direcional efetivamente desaparece. As manchetes, na forma simples de índice de
medo por léxico, parecem contribuir pouco; um tratamento de PLN mais rico é uma
direção promissora. Por fim, a previsão por quantis (Seção~\ref{sec:quantis})
mostrou-se um complemento valioso: em vez de uma estimativa pontual ilusoriamente
precisa, ela comunica a incerteza, sendo essa mais útil para decisão sob risco.

% ===================================================================
\section{Conclusão}
\label{sec:conclusao}

Este trabalho tratou a previsão do preço do ouro como regressão multi-horizonte
de log-retornos e comparou, sob um protocolo único e sem vazamento, cinco
famílias de modelos. Os principais achados são: (i) reformular o alvo em
log-retorno e avaliar contra um \emph{random walk} é essencial, visto que métricas na
escala de preço, como o $R^2$, são enganosas; (ii) num \emph{split} de teste
recente, a CNN 1D causal e o XGBoost superam o passeio aleatório em todos os
horizontes, com $\approx 62\%$ de acurácia direcional em $h=30$; e (iii) sob
validação \emph{walk-forward}, o ganho é pequeno e fortemente dependente do
regime, concentrando-se em anos de tendência e desaparecendo em mercados
laterais.

As principais limitações são o sinal intrinsecamente fraco de séries
financeiras, o uso de um único ativo e período, e o tratamento rudimentar do
texto das manchetes. Como trabalhos futuros, destacam-se: substituir o índice de
medo por léxico por modelos de sentimento financeiro ou \emph{embeddings};
incorporar variáveis macroeconômicas exógenas (juros reais, dólar, outros
metais) de forma causal; avaliar arquiteturas de atenção como o \emph{Temporal
Fusion Transformer} \cite{lim2021tft}; e estender a previsão intervalar por
quantis a todos os modelos, com calibração formal dos intervalos.

% ===================================================================
\section*{Contribuições dos autores}
% PREENCHER com os nomes reais; mapeamento por componente do projeto.
\begin{itemize}
  \item \textbf{Paulino:} coleta e \emph{download} dos dados, análise
        exploratória (EDA), \emph{pipeline} de PLN das manchetes (índice de
        medo) e engenharia de atributos.
  \item \textbf{Ivina:} \emph{baseline} de regressão linear e modelo
        XGBoost com otimização bayesiana; análise por ano.
  \item \textbf{Gabriel:} redes recorrentes (GRU e LSTM) e a CNN 1D causal
        (modelo principal), incluindo a variante \emph{walk-forward} por
        quantis.
  \item \textbf{José:} infraestrutura compartilhada (\texttt{src/}:
        janelamento, métricas, visualização), protocolo de validação,
        rastreamento de experimentos e redação do artigo.
\end{itemize}

\bibliographystyle{sbc}
\bibliography{referencias}

\end{document}